\section{Taust ja seotud tööd} \label{taust}

Käesolevas peatükis antakse ülevaade töö tausta moodustavatest valdkondadest. Alapeatükis~\ref{taust:llm} kirjeldatakse suurte keelemudelite arhitektuurset alust ja võimekust. Alapeatükk~\ref{taust:eesti} käsitleb LLM-ide võimekust eesti keeles ning sellega seotud kitsaskohti. Alapeatükis~\ref{taust:akadeemiline} määratletakse akadeemilise teksti kvaliteediomadused, millele süsteem tagasisidet annab. Alapeatükk~\ref{taust:promptimine} esitab promptimise põhimõtted, mis on prototüübi disainile aluseks. Lõpuks vaadeldakse alapeatükis~\ref{taust:seotud} teksti tagasisidestamise olemasolevaid lahendusi ja akadeemilist kirjandust.

\subsection{Suured keelemudelid} \label{taust:llm}

Suur keelemudel on isejuhendatud õppega närvivõrk, mis on treenitud väga suurel loomuliku keele andmestikul nii, et ta suudab etteantud tekstiosa põhjal ennustada järgnevat sõna~\cite{vaswani_attention_2017, brown_language_2020}. Tänapäevased mudelid põhinevad \emph{Transformer}-arhitektuuril~\cite{vaswani_attention_2017}, milles keskne mehhanism on enesetähelepanu (\emph{ingl.\ self-attention}). See mehhanism võimaldab mudelil iga sõna esituses arvesse võtta kõigi eelnevate (ning mõnel juhul ka järgnevate) sõnade konteksti, mis teeb mudeli pikkade sõltuvuste tuvastamisel oluliselt võimekamaks kui varem domineerinud rekurrentsed närvivõrgud~\cite{hochreiter_lstm_1997}.

Mudeli treenimine jaguneb tüüpiliselt kaheks etapiks: eelkoolitus (\emph{ingl.\ pre-training}) suurtel struktureerimata tekstikorpustel ning järelhäälestus (\emph{ingl.\ fine-tuning}) konkreetseteks ülesanneteks. Juhiseid järgivate mudelite (nagu GPT, Claude, Gemini) puhul lisandub veel inimtagasiside abil õppimine (\emph{ingl.\ reinforcement learning from human feedback}, RLHF)~\cite{ouyang_instructgpt_2022}, mille käigus mudel kohandatakse vastama loomuliku keele juhistele turvaliselt ja kasulikult.

Treeningu mahukus on väga kiiresti kasvanud. Nii GPT-3 (175 miljardit parameetrit, 2020)~\cite{brown_language_2020} kui ka selle järglased on näidanud nähtust, mida on hakatud nimetama \emph{esiletõusvateks võimeteks} (\emph{ingl.\ emergent abilities}): mudelid lahendavad ülesandeid, mille jaoks neid spetsiaalselt ei treenitud, kui parameetrite arv ja andmemaht ületavad teatud piiri~\cite{wei_emergent_2022}. Just see omadus võimaldab kasutada üht ja sama mudelit nii teksti kokkuvõtmiseks, küsimustele vastamiseks kui ka käesolevas töös käsitletud tagasiside andmiseks.

LLM-ide kasutamisega kaasnevad ka tuntud puudujäägid. Mudelid võivad esitada faktiliselt ekslikke väiteid suure enesekindlusega (\emph{ingl.\ hallucination})~\cite{ji_hallucination_2023}, mille olulisust akadeemilise teksti kontekstis ei saa alahinnata. Lisaks võivad mudelid jäljendada treeningandmetes esinevaid süstemaatilisi kallutusi ning ühe konkreetse mudeli väljund ei pruugi olla deterministlik: sama prompt võib eri kordadel anda erinevaid tulemusi.

\subsection{LLM-id eesti keeles} \label{taust:eesti}

Eesti keel on inglise keelega võrreldes vähese ressursiga keel: kuigi tegu ei ole maailma kõige väiksema kõnelejaskonnaga keelega, on eestikeelse digitaalse teksti maht oluliselt väiksem ning seetõttu on eesti keele esindatus tüüpilises mitmekeelses treeningkorpuses (näiteks Common Crawl) väga tagasihoidlik. Sellele vaatamata on viimaste põlvkondade üldotstarbelistes mudelites eesti keele kvaliteet märgatavalt paranenud. Tartu Ülikooli loomuliku keele töötluse uurimisrühma (TartuNLP) avalik eesti keele LLM-võrdlusplatvorm (baromeeter)~\cite{tartunlp_baromeeter_2025} võrdleb nii avatud kui ka kommertsiaalsete mudelite jõudlust eestikeelsetel ülesannetel ning näitab, et tipptasemel mitmekeelsed mudelid (sh Claude 4.x ja GPT-5.5) kuuluvad eesti keele kvaliteedi poolest tippu, kuigi jäävad üldjuhul alla võrreldavale inglise keele tasemele.

Lisaks üldotstarbelistele LLM-idele on Eestis arendatud spetsiaalseid eesti keele mudeleid. EstBERT~\cite{tanvir_estbert_2021} on BERT-arhitektuuril põhinev mudel, mis on eelkoolitatud eestikeelsel korpusel ning sobib selliste klassifitseerimisülesannete jaoks nagu sentiment- ja teemaanalüüs. EstBERT-i kasutamine teksti genereerimist nõudvate ülesannete jaoks on aga piiratud, kuna mudel on arhitektuurselt enkooderitüüpi. Eestikeelseid generatiivseid mudeleid on hakatud aktiivselt arendama hiljutiste avatud lähtekoodiga mudelite (näiteks Llama-2 ja Llama-3) baasil -- nii on TartuNLP koostanud Llammas-mudeli, mille treenimine ja hindamine on dokumenteeritud Kuulmetsa jt töödes~\cite{kuulmets_llammas_2024, kuulmets_finnougri_2025} --, kuid nende ulatus jääb tippkommertsmudelite omale alla.

Praktilises rakenduses tähendab see, et eestikeelse akadeemilise teksti analüüsiks on otstarbekam kasutada tippkvaliteediga mitmekeelseid mudeleid (Claude 4.7 Opus, GPT-5.5) kui spetsialiseeritud avatud mudeleid, kui prioriteediks on kvaliteet, mitte kohaliku andmetöötluse nõue.

\subsection{Akadeemilise teksti kvaliteediomadused} \label{taust:akadeemiline}

Akadeemilise teksti kvaliteet on mitmemõõtmeline. Burstein jt~\cite{burstein_aes_2013} on automaatse esseede hindamise (\emph{ingl.\ automated essay scoring}, AES) ülevaates eristanud vähemalt järgmisi mõõtmeid: sisuline asjakohasus, organiseeritus, sõnavara, stiili sobivus ning õigekirja- ja vormistusvead. Käesolevas töös on rõhk neljal omadusel, mis on Tartu Ülikooli arvutiteaduse instituudi nõuete~\cite{ut_juhend_2024} kontekstis kõige otsesemalt hinnatavad ja mille puhul automaatne tagasiside on praktilise väärtusega.

\subsubsection{Struktuur}

Lõputöö struktuur peab olema lugejale loogiliselt järgitav. Iga peatükk algab sissejuhatava lõiguga, alampeatükid peavad olema teemaga seotud, peatüki sees on hea kasutada siduvaid lauseid eri lõikude vahel ning peatüki lõpus võiks olla sünteesilõik. Peatükk ei peaks algama ega lõppema joonise, tabeli ega loeteluga ilma teksti vaheta~\cite{ut_juhend_2024}. Strukturaalsed probleemid hõlmavad seega liiga lühikesi peatükke, sissejuhatava lõigu puudumist, ebaloogilist alampeatükkide järjekorda ja sujumatut lõikudevahelist üleminekut.

\subsubsection{Akadeemiline stiil}

Akadeemiline stiil eestikeelses lõputöös järgib eesti keele head tava ja teadusliku kirjutamise norme. Tüüpilised stiilivead on järgmised: kõnekeelsete väljendite kasutamine (näiteks \emph{tüüp}, \emph{asi}, \emph{teha}); umbisikuline tegumood seal, kus seda ei nõuta, või vastupidine viga; mina-vormi liigne kasutamine; liiga pikad ja keerulised laused; kvalifitseerimata üldistused (\emph{kõik teavad, et\dots}); inglise keele otsetõlke ilmingud (näiteks \emph{see on tõsi, et\dots} kui konstruktsiooni \emph{it is true that\dots} otsetõlge). Hyland~\cite{hyland_metadiscourse_2005} on rõhutanud, et akadeemiline stiil ei ole pelgalt sõnavara valik, vaid ka selge metadiskursus, mis juhib lugejat läbi teksti.

\subsubsection{Terminoloogia järjepidevus}

Tehnilises informaatikatekstis on terminoloogiavalik eriti oluline. Sama mõiste kohta tuleks kasutada läbivalt sama terminit (näiteks \emph{andmebaas}, mitte vaheldumisi \emph{andmebaas} ja \emph{baas}; või \emph{lõim}, mitte vaheldumisi \emph{lõim} ja \emph{niit}). Tihti tuleb teha valik kahe võrdselt levinud variandi vahel, näiteks \emph{kasutaja\-liides} ja \emph{kasutaja\-pind}, \emph{järjend} ja \emph{loend}, \emph{kompileerimine} ja \emph{tõlkimine}. Eesti Keele Instituudi terminisõnastik\footnote{\url{https://termin.eki.ee/}} ja Tartu Ülikooli arvutiteaduse instituudi sõnastikud annavad osa juhiseid, kuid sageli peab autor tegema iseseisva valiku ning seda järjekindlalt rakendama. Lisaks kerkib küsimus, millal tuua sisse ingliskeelne termin (\emph{kursiivis sulgudes}) ja millal seda enam mitte korrata.

\subsubsection{Viitamisvajadus}

Iga väide, mis ei ole autori enda originaalne mõte ega üldteadmine, vajab viidet allikale. Vastasel juhul on tegu akadeemilise petturlusega. Levinud probleemid on järgmised: viiteta esitatud arvulised väited (\emph{75~\% kasutajatest eelistab\dots}); konkreetne metoodika või algoritm, mille puhul viidet pole esitatud; teiste autorite definitsioonide refereerimine ilma viiteta. LLM-i ülesanne ei ole \emph{ise} viiteid leida, vaid anda märku, kus viide on tõenäoliselt vajalik; lõplik otsus ja kontroll jäävad autorile.

\subsection{Promptimine ja struktureeritud viibad} \label{taust:promptimine}

LLM-i käitumist juhitakse \emph{prompti} ehk viiba abil – see on tekstiline juhis, mille mudel saab koos analüüsitava sisendiga. Viiba kvaliteet mõjutab mudeli väljundit oluliselt. Promptimisstrateegiate ülevaates~\cite{liu_prompting_2023} eristatakse muu hulgas:

\begin{itemize}
    \item \textbf{Nullhetk-promptimine} (\emph{ingl.\ zero-shot prompting}) – mudelile antakse ainult ülesanne ja sisend, näideteta.
    \item \textbf{Vähehetk-promptimine} (\emph{ingl.\ few-shot prompting}) – ülesandele lisatakse mõned näited oodatava sisendi-väljundi paari kohta.
    \item \textbf{Mõtteahel-promptimine} (\emph{ingl.\ chain-of-thought prompting}) – mudelilt küsitakse selgesõnaliselt vahesammude põhjendust enne lõplikku vastust~\cite{wei_chain_2022}.
    \item \textbf{Struktureeritud väljundi nõue} – mudelilt nõutakse väljundit kindlal kujul, näiteks JSON-i, mis lihtsustab automaatset järeltöötlust~\cite{geng_structured_2025}.
\end{itemize}

Käesolevas töös testitakse kaht promptimisstrateegiat. Mõlemad nõuavad mudelilt sama JSON-skeemiga väljundit, et automaatne järeltöötlus oleks võrreldav. \emph{Üldine} (\emph{baseline}) prompt esitab mudelile vaid rolli ja minimaalse JSON-skeemi (lisaks neljale tagasisidekategooriale lubatud ka \texttt{MUU}), kuid \emph{ei sisalda} kategooriate definitsioone, näiteid ega täpsustavaid reegleid. \emph{Struktureeritud} prompt määrab täpselt neli kategooriat, esitab iga kategooria kohta näiteid ja vastunäiteid, sisaldab mõtteahela elementi (mudel peab iga tuvastatud probleemi jaoks põhjendama, miks ta seda probleemiks peab) ning sõnastab konservatiivsusele suunavaid reegleid. Hüpotees on, et struktureeritud prompt annab täpsemat ja võrreldavamat väljundit.

\subsection{Seotud tööd} \label{taust:seotud}

Automaatse teksti tagasisidestamise valdkonnal on pikk ajalugu. Klassikalised süsteemid, nagu Criterion~\cite{burstein_criterion_2004} ja e-rater~\cite{attali_erater_2006}, on alates 2000. aastatest võimaldanud esseetasemel tagasisidet, kuid on välja töötatud peamiselt inglise keele jaoks. Grammarly~\cite{grammarly_2024} on tänapäeval üks enim kasutatud kommertsiaalseid stiili- ja grammatikakontrolli vahendeid, mille eesti keele tugi on 2026.~aasta seisuga endiselt piiratud.

Teadusartiklite kvaliteedi automaatset tagasisidestamist on uuritud süsteemides, nagu SciScore~\cite{menke_sciscore_2020} ja Paper Robot~\cite{wang_paperrobot_2019}, kuid need keskenduvad metoodika reprodutseeritavusele ja struktureeritud kontrollnimekirjadele, mitte üldisele stiilile ja terminoloogiale. LLM-ide tulekuga on hakanud ilmuma tööd, mis kasutavad GPT-4 ja Claude'i akadeemiliste kirjutiste tagasisidestamiseks~\cite{liang_chatgpt_2024, kasneci_chatgpt_2023}. Liang jt~\cite{liang_chatgpt_2024} näitavad, et GPT-4 teadusartiklitele antud tagasiside kattub umbes 30--39~\% ulatuses inimretsensentide tagasisidega; see protsent on madalam kui inimretsensentide omavaheline kattuvus, kuid sellele lähedane.

Eestikeelse akadeemilise teksti automaatset tagasisidestamist on uuritud oluliselt vähem. Käesoleva töö raames ei õnnestunud tuvastada avaldatud uurimusi, mis käsitleksid spetsiaalselt eestikeelsete informaatika lõputööde automaatset tagasisidestamist generatiivsete LLM-ide abil. Käesolev töö täidab seda lünka kahel viisil: esiteks rakendab süsteemi spetsiaalselt informaatika lõputöödele, kus tehniline terminoloogia mängib suurt rolli, ja teiseks kasutab generatiivseid LLM-e koos struktureeritud promptiga.

Promptimise mõju kvaliteedile on samuti uuritud. Schulhoff jt~\cite{schulhoff_prompt_report_2024} on promptimistehnikaid süstemaatiliselt katalogiseerinud ning näitavad, et struktureeritud promptid annavad spetsiifilistel ülesannetel tüüpiliselt suuremat täpsust. Käesolev töö rakendab seda eestikeelse akadeemilise teksti analüüsile.

Käesoleva töö \textbf{teaduslik panus} on seega: (1) struktureeritud LLM-prompti disain eestikeelse informaatikateksti tagasisidestamiseks neljas kategoorias, (2) 20-katkendilise sünteetilise testkogu ja kuldstandardi koostamine koos hindamismetoodikaga ning (3) üldise ja struktureeritud prompti empiiriline võrdlus kahe tippmudeli põhjal, mille tulemused näitavad struktureeritud prompti järjepidevat eelist.
